\subsection{\texorpdfstring{第~\ref{sec:differentiable-manifolds}~节}{可微流形}}\label{subsec:solutions-differentiable-structures}

\begin{description}[font=\normalfont,leftmargin=2em,labelwidth=1.5em]
\RGsolution{ex:prove-inequivalence-of-line-atlases}{prove-inequivalence-of-line-atlases}
$(\varphi_2^{-1}\circ\varphi_1)(x)=x^{1/3}$ 在原点不可微. 由 Exercise~\ref{ex:exercise-differentiable-structures}(\ref{ex:compare-atlases-by-union}), 由于 $\{(\mathbb{R},\varphi_1),(\mathbb{R},\varphi_2)\}$ 不是图册, 这两个图册不等价.

\RGsolution{ex:stereographic-sphere-atlas}{stereographic-sphere-atlas}
把点 $p\in S^n\subset\mathbb{R}^{n+1}$ 写成 $p=(x,x^{n+1})$, 其中 $x\in\mathbb{R}^n$, $x^{n+1}\in\mathbb{R}$. 经过 $N$ 和 $p$ 的直线的参数方程为
\[
  c(t)=(tx,1+t(x^{n+1}-1)),
\]
它在 $t=\frac{1}{1-x^{n+1}}$ 时与超平面 $x^{n+1}=0$ 相交. 因此,
\[
  y=\pi_N(p)=\frac{x}{1-x^{n+1}}.
\]
对等式两边取范数的平方, 得
\[
  \lVert y\rVert^2
  =\frac{\lVert x\rVert^2}{(1-x^{n+1})^2}
  =\frac{1-(x^{n+1})^2}{(1-x^{n+1})^2}
  =\frac{1+x^{n+1}}{1-x^{n+1}}
\]
(这里用了 $\lVert p\rVert^2=\lVert x\rVert^2+(x^{n+1})^2=1$); 等价地,
\[
  x^{n+1}=\frac{\lVert y\rVert^2-1}{\lVert y\rVert^2+1}.
\]
所以 $x$ 与 $y$ 的关系可写成
\[
  x=\frac{2y}{1+\lVert y\rVert^2},
\]
从而
\[
  p=\pi_N^{-1}(y)
  =\left(\frac{2y}{1+\lVert y\rVert^2},\frac{\lVert y\rVert^2-1}{\lVert y\rVert^2+1}\right).
\]
类似地,
\[
  y=\pi_S(p)=\frac{x}{1+x^{n+1}}
\]
且
\[
  p=\pi_S^{-1}(y)
  =\left(\frac{2y}{1+\lVert y\rVert^2},\frac{1-\lVert y\rVert^2}{1+\lVert y\rVert^2}\right).
\]
因此, 把 $\pi_N(S^n\setminus\{N,S\})=\mathbb{R}^n\setminus\{0\}$ 映到 $\pi_S(S^n\setminus\{N,S\})=\mathbb{R}^n\setminus\{0\}$ 的映射 $\pi_S\circ\pi_N^{-1}$ 由
\[
  \pi_S\circ\pi_N^{-1}(y)=\frac{y}{\lVert y\rVert^2}
\]
给出, 因而在其定义域上可微. 同样,
\[
  \pi_N\circ\pi_S^{-1}(y)=\frac{y}{\lVert y\rVert^2}
\]
也在其定义域上可微. 因此 $\{(\mathbb{R}^n,\pi_N^{-1}),(\mathbb{R}^n,\pi_S^{-1})\}$ 是 $S^n$ 的一个图册.

要证明这个图册与 Example~\ref{ex:examples-of-differentiable-structures}(\ref{ex:sphere-atlas-by-coordinate-projections}) 中的图册等价, 需要计算 $i=1,\ldots,n+1$ 时的 $\pi_N\circ\varphi_i^+$, $\pi_N\circ\varphi_i^-$, $\pi_S\circ\varphi_i^+$ 和 $\pi_S\circ\varphi_i^-$ 及其逆. 本质上有两种情形, 分别对应于 $i=n+1$ 和 $i\ne n+1$. 例如,
\[
  \pi_N\circ\varphi_{n+1}^+(x)
  =\pi_N(x,g(x))
  =\frac{x}{1-g(x)}
  =\frac{x}{1-(1-\lVert x\rVert^2)^{\frac12}}
\]
在其定义域 $U\setminus\{0\}$ 上可微. 另一种情形类似.

\RGsolution{ex:projective-space-atlas-from-lines}{projective-space-atlas-from-lines}
\begin{enumerate}[label=(\alph*),ref=\alph*,leftmargin=*]
\item 显然 $\mathbb{R}P^n=\bigcup_{i=1}^{n+1} V_i$. 此外, 若 $\varphi_i(x^1,\ldots,x^n)\in V_j$, 则
\begin{align*}
  \varphi_i(x^1,\ldots,x^n)
  &=[x^1,\ldots,x^j,\ldots,x^{i-1},1,x^i,\ldots,x^n]\\
  &=\left[\frac{x^1}{x^j},\ldots,1,\ldots,\frac{x^{i-1}}{x^j},\frac{1}{x^j},\frac{x^i}{x^j},\ldots,\frac{x^n}{x^j}\right],
\end{align*}
因而
\[
  \varphi_j^{-1}\circ\varphi_i(x^1,\ldots,x^n)
  =\left(\frac{x^1}{x^j},\ldots,\frac{x^{i-1}}{x^j},\frac{1}{x^j},\frac{x^i}{x^j},\ldots,\frac{x^n}{x^j}\right)
\]
在其定义域上可微.
\item 商空间 $(\mathbb{R}^{n+1}\setminus\{0\})/{\sim}$ 和 $S^n/{\sim}$ 之间有一个双射, 因为过原点的每条直线都在一对对径点处与 $S^n$ 相交. 两个拓扑也相同, 因为 $\mathbb{R}^{n+1}\setminus\{0\}$ 中由过原点的直线形成的开集与 $S^n$ 的交集是一个开集. 要检验这两个图册等价, 需要检验当 $i,j=1,\ldots,n+1$ 时, 映射 $\varphi_j^{-1}\circ(\pi\circ\varphi_i^\pm)$ 及其逆可微. 例如,
\begin{align*}
  \varphi_{n+1}^{-1}\circ(\pi\circ\varphi_{n+1}^+)(x)
  &=\varphi_{n+1}^{-1}(\pi(x,g(x)))
  =\varphi_{n+1}^{-1}([x,g(x)])\\
  &=\varphi_{n+1}^{-1}\left(\left[\frac{x}{g(x)},1\right]\right)
  =\frac{x}{g(x)},
\end{align*}
这在其定义域 $U$ 上可微. 其他情形类似.
\end{enumerate}
\end{description}
